\section{Period-coordinate Dolbeault classes and the Bott--Chern algebra}
\label{sec:bcpr-product-structure}

\subsection{Fixed input and conventions}

Let $f\colon X\to B=\mathbb P^1$ be the compact complex threefold of
\cref{thm:candidate-global-geometry}.  We use
$d=\partial+\bar\partial$ and, following
\cite[\S1.a]{Schweitzer2007},
\begin{equation}
 H^{p,q}_{BC}(X)=
 \frac{\ker\partial\cap\ker\bar\partial
       \subset\mathcal A^{p,q}(X)}
      {\operatorname{im}(\partial\bar\partial)}.
\label{eq:bcpr-bc-definition}
\end{equation}
The natural map to Dolbeault cohomology is denoted
\begin{equation}
 \pi_{\bar\partial}^{p,q}\colon H^{p,q}_{BC}(X)
 \longrightarrow H^{p,q}_{\bar\partial}(X).
\label{eq:bcpr-dolbeault-map}
\end{equation}
Exterior multiplication makes the direct sum of the Bott--Chern groups a
bigraded algebra, and complex conjugation interchanges bidegrees $(p,q)$
and $(q,p)$ \cite[\S1.a]{Schweitzer2007}.

By \cref{cor:exact-hodge-diamond}, the two nonzero rows of the
Fr\"olicher differential are the exact complexes
\begin{equation}
 0\longrightarrow H^{0,1}_{\bar\partial}(X)
 \xrightarrow{d_1}H^{1,1}_{\bar\partial}(X)
 \xrightarrow{d_1}H^{2,1}_{\bar\partial}(X)
 \longrightarrow0,
\label{eq:bcpr-hodge-row-one}
\end{equation}
\begin{equation}
 0\longrightarrow H^{1,2}_{\bar\partial}(X)
 \xrightarrow{d_1}H^{2,2}_{\bar\partial}(X)
 \xrightarrow{d_1}H^{3,2}_{\bar\partial}(X)
 \longrightarrow0,
\label{eq:bcpr-hodge-row-two}
\end{equation}
and each row has dimensions $1,2,1$.  The exact Bott--Chern calculation
of \cref{thm:exact-bott-chern-aeppli}, obtained from the compact-threefold
formulae of \cite{McHugh2019}, gives
\begin{equation}
\begin{array}{c|cccccccc}
(p,q)&(0,0)&(1,1)&(2,1)&(1,2)&(2,2)&(3,2)&(2,3)&(3,3)\\ \hline
h^{p,q}_{BC}&1&2&1&1&2&1&1&1.
\end{array}
\label{eq:bcpr-nonzero-dimensions}
\end{equation}
All other Bott--Chern groups vanish.  By
\cref{cor:candidate-sphere}, $b_k(X)=0$ for $1\leq k\leq5$.

On a marked smooth fibre retain the integral basis
\[
 V=\mathbb Z\langle\gamma,u,w,\delta\rangle
\]
without renaming its element $u$, and put
\begin{equation}
 \eta=u\wedge w+6\gamma\wedge\delta,
 \qquad
 \eta^2=12\gamma\wedge u\wedge w\wedge\delta.
\label{eq:bcpr-eta-square}
\end{equation}
The holomorphic Hodge rows supplied by
\cref{prop:period-input-audit,prop:very-general-torus-audit} are
\begin{equation}
 \sigma_1=6\mu\gamma+\tau u+w,
 \qquad
 \sigma_2=\beta\gamma+\mu u+\delta.
\label{eq:bcpr-period-rows}
\end{equation}
Thus no global Dolbeault or Bott--Chern generator below is denoted by
$u$ or $x$; those letters remain reserved for the printed fibre basis and
for local real linear functions.

\subsection{Exact fibre representatives}

Write
\begin{equation}
 y=\operatorname{Im}\tau,
 \qquad m=\operatorname{Im}\mu,
 \qquad r=\operatorname{Im}\beta,
 \qquad \Delta=r-\frac{6m^2}{y}<0.
\label{eq:bcpr-delta}
\end{equation}
The four marked lattice columns in the fibre coordinate
$\zeta=(\zeta_1,\zeta_2)$ are
\begin{equation}
 (6\mu,\beta),\quad(\tau,\mu),\quad(1,0),\quad(0,1),
\label{eq:bcpr-lattice-columns}
\end{equation}
in the order $\widehat\gamma,\widehat u,\widehat w,
\widehat\delta$.  On the universal cover of a marked fibre define the
real linear functions
\begin{equation}
 x_\gamma=
 \frac{y\,\operatorname{Im}\zeta_2-m\,\operatorname{Im}\zeta_1}
      {y\Delta},
 \qquad
 x_u=
 \frac{r\,\operatorname{Im}\zeta_1-6m\,\operatorname{Im}\zeta_2}
      {y\Delta}.
\label{eq:bcpr-real-linear-functions}
\end{equation}
For the first two columns in \cref{eq:bcpr-lattice-columns}, the numerators
of the changes in $x_\gamma$ are respectively
$yr-6m^2=y\Delta$ and $ym-my=0$, while the numerators of the changes in
$x_u$ are $6rm-6mr=0$ and $ry-6m^2=y\Delta$.  The last two columns are
real and change neither imaginary part.  Hence
\[
 \bigl(\Delta_{\widehat\gamma}x_\gamma,
       \Delta_{\widehat u}x_\gamma,
       \Delta_{\widehat w}x_\gamma,
       \Delta_{\widehat\delta}x_\gamma\bigr)=(1,0,0,0),
\]
\begin{equation}
 \bigl(\Delta_{\widehat\gamma}x_u,
       \Delta_{\widehat u}x_u,
       \Delta_{\widehat w}x_u,
       \Delta_{\widehat\delta}x_u\bigr)=(0,1,0,0).
\label{eq:bcpr-linear-function-increments}
\end{equation}
Their differentials, although not the functions themselves, therefore
descend to the fibre.  Their $(0,1)$ parts are
\begin{equation}
 \Gamma:=\bar\partial_Fx_\gamma
 =\frac{i}{2\Delta}
 \left(d\bar\zeta_2-\frac{m}{y}d\bar\zeta_1\right),
\label{eq:bcpr-gamma-form}
\end{equation}
\begin{equation}
 U:=\bar\partial_Fx_u
 =\frac{i}{2y\Delta}
 \left(r\,d\bar\zeta_1-6m\,d\bar\zeta_2\right).
\label{eq:bcpr-u-form}
\end{equation}
They represent respectively the $(0,1)$ images of the integral classes
$\gamma$ and $u$.

The identity
\begin{equation}
 \eta=u\wedge\sigma_1+6\gamma\wedge\sigma_2
     =-\sigma_1\wedge u-6\sigma_2\wedge\gamma
\label{eq:bcpr-eta-decomposition}
\end{equation}
holds in $\bigwedge^2V_{\mathbb C}$.  Since $\eta$ is of type $(1,1)$
on every marked fibre, \cref{eq:bcpr-eta-decomposition} gives its exact
Dolbeault representative
\begin{equation}
 Q_F=-d\zeta_1\wedge U-6d\zeta_2\wedge\Gamma.
\label{eq:bcpr-q-fibre}
\end{equation}
Equivalently, with the Appell--Humbert convention
$\eta=(2i)^{-1}\sum(H_\eta)_{ab}\sigma_a\wedge\bar\sigma_b$, the matrix
calculation in the proof of \cref{prop:very-general-torus-audit} gives
\begin{equation}
 H_\eta=\frac1{y\Delta}
 \begin{pmatrix}r&-6m\\-6m&6y\end{pmatrix},
\label{eq:bcpr-appell-humbert-matrix}
\end{equation}
and therefore
\begin{equation}
 Q_F=\frac1{2iy\Delta}\Bigl(
 r\,d\zeta_1\wedge d\bar\zeta_1
 -6m\,d\zeta_1\wedge d\bar\zeta_2
 -6m\,d\zeta_2\wedge d\bar\zeta_1
 +6y\,d\zeta_2\wedge d\bar\zeta_2\Bigr).
\label{eq:bcpr-q-fibre-expanded}
\end{equation}
Direct exterior multiplication, with holomorphic differentials ordered
before antiholomorphic differentials, gives
\begin{equation}
 Q_F\wedge\Gamma
 =\frac1{4y\Delta}
 d\zeta_1\wedge d\bar\zeta_1\wedge d\bar\zeta_2\neq0,
\label{eq:bcpr-q-gamma-product}
\end{equation}
\begin{equation}
 Q_F^2
 =\frac3{y\Delta}
 d\zeta_1\wedge d\zeta_2\wedge
 d\bar\zeta_1\wedge d\bar\zeta_2\neq0.
\label{eq:bcpr-q-square}
\end{equation}
Indeed,
$\det H_\eta=(6ry-36m^2)/(y^2\Delta^2)=6/(y\Delta)$,
and the coefficient in \cref{eq:bcpr-q-square} is
$\tfrac12\det H_\eta$ in the displayed ordering.  The identities
\Cref{eq:bcpr-q-gamma-product,eq:bcpr-q-square} will be used only in
fibre cohomology.

\subsection{Extension of the two Dolbeault inputs}

\begin{lemma}[the degree-one input]
\label{lem:bcpr-degree-one-input}
\Verified\ There is a distinguished class
$\mathfrak g\in H^{0,1}_{\bar\partial}(X)$ whose restriction to every
smooth marked fibre is represented by $\Gamma$.
\end{lemma}

\begin{proof}
By \cref{thm:structure-direct-images-audit}, the monodromy-invariant
integral class $\gamma$ maps to the nowhere-zero holomorphic section
$\gamma^{0,1}$ of
$R^1f_*\mathcal O_X\simeq\mathcal O\oplus\mathcal O(-1)$.  Coherent
Leray gives
\[
 H^1(X,\mathcal O_X)=H^0(B,R^1f_*\mathcal O_X)
 =\mathbb C\gamma^{0,1}.
\]
Let $\mathfrak g$ be the class mapping to $\gamma^{0,1}$.
\Cref{eq:bcpr-gamma-form} is the exact period-coordinate
representative of its fibre restriction.  This is the distinguished
coordinate of \cref{cor:picard-line-audit} before any analytic rescaling.
\end{proof}

\begin{lemma}[extension of the invariant $(1,1)$ class]
\label{lem:bcpr-eta-extension}
\Verified\ The fibre class $\eta$ extends to a nonzero holomorphic section
\begin{equation}
 \boldsymbol\eta\in
 H^0\bigl(B,R^1f_*\Omega_X^1\bigr).
\label{eq:bcpr-eta-global-section}
\end{equation}
Consequently there is a class
$\mathfrak q\in H^{1,1}_{\bar\partial}(X)$ whose coherent-Leray edge
value is $\boldsymbol\eta$.  For every smooth fibre inclusion
$i_b\colon F_b\to X$,
\begin{equation}
 i_b^*\mathfrak q=[Q_F]_{\bar\partial}
 \in H^{1,1}_{\bar\partial}(F_b).
\label{eq:bcpr-q-restriction}
\end{equation}
\end{lemma}

\begin{proof}
Over the smooth marked locus, $\eta$ is flat, monodromy invariant, and of
type $(1,1)$ identically.  Equation
\Cref{eq:bcpr-eta-decomposition} identifies the corresponding holomorphic
section of the relative group as
\begin{equation}
 -\sigma_1\wedge u^{0,1}-6\sigma_2\wedge\gamma^{0,1}.
\label{eq:bcpr-relative-eta-section}
\end{equation}

At a finite filling write $t=s^M$, with $(M,a)=(3,2)$ or $(4,1)$.
The covering family is smooth at $s=0$, and
\cref{eq:bcpr-eta-decomposition} holds identically there, so
\cref{eq:bcpr-relative-eta-section} remains holomorphic.  It is invariant
under the cyclic action.  In equivariant frames $e_0,e_1$ for the relative
cotangent and $b_0,b_1$ for $R^1\mathcal O$, the actual invariant
$\mathbb C\{t\}$-basis from
\cref{prop:log-relative-direct-images} is
\begin{equation}
 e_0b_0,\quad s^ae_0b_1,\quad
 s^{M-a}e_1b_0,\quad e_1b_1.
\label{eq:bcpr-finite-invariant-basis}
\end{equation}
Exactness of cyclic invariants in characteristic zero therefore makes
\cref{eq:bcpr-relative-eta-section} a holomorphic descended section at
both finite points.

At the cusp, $\gamma$ and $u$ extend as sections of $R^1f_*\mathbb C$,
and their $(0,1)$ images form a basis of $H^1(W,\mathcal O_W)$ by
\cref{thm:structure-direct-images-audit}.  On the dense torus,
\[
 \sigma_i=d\zeta_i=(2\pi i)^{-1}d\log x_i.
\]
The classes of these logarithmic forms extend in the relative logarithmic
sheaf as
\begin{equation}
 \mathcal E|_{N_0}
 =\mathcal O_{N_0}[d\log x_1]\oplus
  \mathcal O_{N_0}[d\log x_2]
\label{eq:bcpr-cusp-log-basis}
\end{equation}
by \cref{prop:log-relative-direct-images}.  A cusp quotient shear changes
$d\log x_i$ only by a holomorphic multiple of $d\log t$, which is zero in
the relative logarithmic sheaf.  Thus
\cref{eq:bcpr-relative-eta-section} also extends holomorphically through
the cusp.  This concerns classes in the relative logarithmic sheaf; it
does not assert that the $d\log x_i$ are global ordinary one-forms on $X$.

The Gauss--Manin/Kodaira--Spencer map
\[
 \kappa_1\colon R^1f_*\mathcal E\longrightarrow
 \Omega_B^1(\log\Sigma)\otimes R^2f_*\mathcal O_X
\]
annihilates \cref{eq:bcpr-relative-eta-section} on the smooth locus: a
flat class which remains of type $(1,1)$ has zero $(0,2)$ Higgs component.
Both the section and $\kappa_1$ extend holomorphically, so the identity
theorem gives $\kappa_1(\boldsymbol\eta)=0$ on $B$.  The ordinary-cotangent
calculation in
\cref{prop:log-ks-divisor,prop:absolute-cotangent-direct-images,thm:h12-scalar}
identifies
\begin{equation}
 R^1f_*\Omega_X^1=\ker\kappa_1=:\mathcal K,
 \qquad
 \mathcal K\simeq\mathcal O\oplus\mathcal O(-1)^{\oplus2}.
\label{eq:bcpr-cotangent-kernel}
\end{equation}
This proves \cref{eq:bcpr-eta-global-section}.  Its value on every smooth
fibre is nonzero in Dolbeault cohomology, not merely as a differential form:
by \cref{eq:bcpr-eta-square} its square is the integral top class
$12\gamma\wedge u\wedge w\wedge\delta$, and therefore
\begin{equation}
 \int_{(F_b,\mathrm{or}_{\rm mark})}Q_F^2=12,
 \qquad
 \int_{(F_b,\mathrm{or}_{\mathbb C})}Q_F^2=-12
\qquad(b\in B^\circ),
\label{eq:bcpr-q-square-integral}
\end{equation}
where $\mathrm{or}_{\rm mark}$ is the orientation determined by
$(\gamma,u,w,\delta)$ and $\mathrm{or}_{\mathbb C}$ is the complex
orientation.  They are opposite because the determinant of the four marked
lattice columns in real coordinates is $y\Delta<0$.  Thus
$[Q_F]_{\bar\partial}\ne0$, so the holomorphic
section $\boldsymbol\eta$ is nonzero.  Explicitly, if
$Q_F=\bar\partial_F\alpha$ for a $(1,0)$-form $\alpha$, then
$Q_F^2=\bar\partial_F(\alpha\wedge Q_F)$; integration and Stokes' theorem
would contradict \cref{eq:bcpr-q-square-integral}.  Here
$\alpha\wedge Q_F$ has bidegree $(2,1)$, so its $\partial_F$ is zero in
complex dimension two and the displayed $\bar\partial_F$-boundary is also a
$d$-boundary.

Finally, coherent Leray over the curve gives the exact edge sequence
\begin{equation}
 0\longrightarrow H^1(B,f_*\Omega_X^1)
 \longrightarrow H^1(X,\Omega_X^1)
 \longrightarrow H^0(B,R^1f_*\Omega_X^1)
 \longrightarrow0,
\label{eq:bcpr-cotangent-edge}
\end{equation}
where $f_*\Omega_X^1\simeq\mathcal O(-2)$ by
\cref{prop:absolute-cotangent-direct-images}.  Choose $\mathfrak q$ over
$\boldsymbol\eta$.  Base change on a smooth fibre gives
\cref{eq:bcpr-q-restriction}.  Any two choices differ by the
one-dimensional left term in \cref{eq:bcpr-cotangent-edge}.  That term
restricts to zero on every smooth fibre.  The nonzero class
$d_1\mathfrak g$ also restricts to zero, because its restriction is
$[\partial_F\Gamma]_{\bar\partial}=0$; since
$H^1(B,\mathcal O(-2))$ is one-dimensional, the left term is exactly
$\operatorname{im}(d_1\colon H^{0,1}_{\bar\partial}\to
H^{1,1}_{\bar\partial})$.
\end{proof}

\begin{remark}[what is and is not explicit]
\label{rem:bcpr-literal-gauge}
Choose global $\bar\partial$-closed forms
\begin{equation}
 G\in\mathcal A^{0,1}(X),\qquad
 Q\in\mathcal A^{1,1}(X)
\label{eq:bcpr-global-representatives}
\end{equation}
representing $\mathfrak g$ and $\mathfrak q$.
\Cref{eq:bcpr-q-restriction} says that
$i_b^*[Q]_{\bar\partial}$ is the Dolbeault class of the explicit
translation-invariant
form displayed in \cref{eq:bcpr-q-fibre-expanded}; it does not say that an arbitrary
representative $Q$ restricts literally to that form.  A literal global
coefficient formula would require additional $C^\infty$ gauge choices and
horizontal correction forms on the disc--pair-of-pants cover, compatible
with the two affine cyclic conjugacies and the toric cusp.  Those choices
are not part of the canonical data.  No such literal gauge is needed for
the class and product calculations below.
\end{remark}

\subsection{The factor \texorpdfstring{$12$}{12} and the Bott--Chern first Chern class}

Let
\[
 \exp\colon H^1(X,\mathcal O_X)\longrightarrow\operatorname{Pic}(X)
\]
be induced by $z\mapsto e^{2\pi iz}$.  By
\cref{prop:leray-e2},
\begin{equation}
 H^0(B,R^1f_*\mathbb Z)=\mathbb Z(12\gamma).
\label{eq:bcpr-global-integral-line}
\end{equation}
The exponential--Leray comparison gives a line bundle $M$ on $B$ such
that
\begin{equation}
 \exp(12\mathfrak g)=f^*M,
 \qquad c_1(M)=d_2(12\gamma).
\label{eq:bcpr-exponential-transgression}
\end{equation}
For completeness, the first assertion follows because
\cref{eq:bcpr-global-integral-line} says that the restriction of
$12\mathfrak g$ to every fibre lies in the image of integral cohomology
and hence exponentiates to the trivial bundle.  The low-degree Leray
sequence for $\mathcal O^*$ then makes the resulting bundle a pullback.
On a disc--pair-of-pants cover of $B$, choose local integral lifts of
$12\gamma$.  Their differences on the annular overlaps form a Cech
one-cochain, and its Cech coboundary is simultaneously the two-cocycle
defining the Leray transgression and the two-cocycle of logarithms defining
$c_1(M)$.  This proves the second assertion with the same sign convention;
it is the exponential-sequence version of the local-lift calculation in
the proof of \cref{lem:first-transgression}.

Orient the generator $\omega\in H^2(B;\mathbb Z)$ by
\begin{equation}
 \omega=c_1(\mathcal O_B(1)).
\label{eq:bcpr-base-orientation}
\end{equation}
Then \cref{lem:first-transgression} and
\cref{eq:bcpr-exponential-transgression} give
\begin{equation}
 \exp(12\mathfrak g)=f^*\mathcal O_B(1).
\label{eq:bcpr-factor-twelve}
\end{equation}
If the opposite orientation is used for the primitive generator in
\cref{lem:first-transgression}, the corresponding statement is
$\exp(-12\mathfrak g)=f^*\mathcal O_B(1)$.  Thus the canonical integral
data fixes the absolute factor $12$, while its sign is the orientation
choice fixed in \cref{eq:bcpr-base-orientation}.

Put
\begin{equation}
 a=\partial G,
 \qquad a^\dagger=\bar\partial\bar G,
 \qquad
 \mathsf A=[a]_{BC},
 \qquad \mathsf B=[a^\dagger]_{BC}.
\label{eq:bcpr-a-b-definitions}
\end{equation}
Both forms are $d$-closed.  We now compute their Chern-class normalization
without choosing local metric weights.  For $z\in\mathbb C$, equip the
smooth trivial line with a global frame $s$ and the holomorphic structure
\begin{equation}
 \bar\partial_zs=-2\pi i zG\otimes s.
\label{eq:bcpr-holomorphic-structure}
\end{equation}
It is integrable because $\bar\partial G=0$.  If
$G=\bar\partial\phi_i$ on a good cover, then
$e_i=e^{2\pi iz\phi_i}s$ is $\bar\partial_z$-holomorphic and
\[
 e_j=e^{2\pi iz(\phi_j-\phi_i)}e_i.
\]
The holomorphic cocycle $\phi_j-\phi_i$ represents $\mathfrak g$ under
the Cech--Dolbeault comparison, so
\cref{eq:bcpr-holomorphic-structure} is precisely the line bundle
$\exp(z\mathfrak g)$ for the displayed exponential convention.

Give $s$ unit length.  The unit-metric Chern connection is
\begin{equation}
 \nabla_z=d-2\pi i\bigl(zG+\bar z\,\bar G\bigr).
\label{eq:bcpr-unit-chern-connection}
\end{equation}
Indeed, its $(0,1)$ part is
\cref{eq:bcpr-holomorphic-structure}, and its connection one-form is
purely imaginary, so it is metric compatible; uniqueness identifies it
with the Chern connection.  Since $\bar\partial G=0$ and
$\partial\bar G=0$, its curvature is the global $(1,1)$ form
\begin{equation}
 F_{\nabla_z}
 =-2\pi i\bigl(z\,\partial G+
                 \bar z\,\bar\partial\bar G\bigr)
 =-2\pi i\bigl(za+\bar z\,a^\dagger\bigr).
\label{eq:bcpr-unit-chern-curvature}
\end{equation}
With the Bott--Chern convention
\begin{equation}
 c_{1,BC}(L,h)=
 \left[\frac{i}{2\pi}F_{\nabla^h}\right]_{BC}
\label{eq:bcpr-chern-convention}
\end{equation}
of \cite[\S8.a]{Schweitzer2007},
\cref{eq:bcpr-unit-chern-curvature} gives, for every $z\in\mathbb C$,
\begin{equation}
 c_{1,BC}\bigl(\exp(z\mathfrak g)\bigr)
 =z\mathsf A+\bar z\,\mathsf B.
\label{eq:bcpr-chern-exp-identity}
\end{equation}
Consequently, if $\varepsilon=1$ for the orientation fixed by
\cref{eq:bcpr-base-orientation} and $\varepsilon=-1$ for its opposite,
then
\begin{equation}
 c_{1,BC}\bigl(f^*\mathcal O_B(1)\bigr)
 =12\varepsilon\bigl(\mathsf A+\mathsf B\bigr).
\label{eq:bcpr-c1-pullback}
\end{equation}
This is the convention-consistent form of the factor-$12$ identity; no
analytic rescaling of the Picard coordinate changes its integral content.

\subsection{All Bott--Chern generators}

Using \cref{eq:bcpr-global-representatives}, define the following
$d$-closed forms and their classes:
\begin{equation}
\begin{array}{c|c|c}
\text{class}&\text{representative}&\text{bidegree}\\ \hline
\mathsf A&a=\partial G&(1,1)\\
\mathsf B&a^\dagger=\bar\partial\bar G&(1,1)\\
\mathsf C&c=\partial Q&(2,1)\\
\overline{\mathsf C}&c^\dagger=\bar\partial\bar Q&(1,2)\\
\mathsf D&d_0=\partial(G\wedge Q)&(2,2)\\
\overline{\mathsf D}&d_0^\dagger=
 \bar\partial(\bar G\wedge\bar Q)&(2,2)\\
\mathsf E&e=\partial(Q^2)&(3,2)\\
\overline{\mathsf E}&e^\dagger=\bar\partial(\bar Q^2)&(2,3)\\
\mathsf T&\nu,\quad\int_X\nu=1&(3,3).
\end{array}
\label{eq:bcpr-generators}
\end{equation}
Here $\nu$ is any smooth positive top-degree form divided by its integral.
The chosen $Q$ cannot be real: its fibre class is real, but the proof below
shows that $\partial Q$ is nonzero globally, whereas a real
$\bar\partial$-closed $Q$ would also be $\partial$-closed.

\begin{proposition}[nonvanishing and bases]
\label{prop:bcpr-generators-are-bases}
\Verified\ Together with the unit $\mathbf1\in H^{0,0}_{BC}(X)$, the entries
in \cref{eq:bcpr-generators} give bases of all the nonzero Bott--Chern groups
in \cref{eq:bcpr-nonzero-dimensions}:
$\mathsf A,\mathsf B$ form the basis in bidegree $(1,1)$ and
$\mathsf D,\overline{\mathsf D}$ form the basis in bidegree $(2,2)$;
each other listed bidegree is one-dimensional and is generated by its
displayed class.
\end{proposition}

\begin{proof}
Because $X$ is compact and connected, a smooth function annihilated by both
$\partial$ and $\bar\partial$ is constant.  There is no
$\partial\bar\partial$ image in bidegree $(0,0)$, so
$\mathbf1:=[1]_{BC}$ generates
$H^{0,0}_{BC}(X)=\mathbb C\mathbf1$.

The Dolbeault image of $\mathsf A$ is
$d_1\mathfrak g\neq0$ by injectivity in
\cref{eq:bcpr-hodge-row-one}.  Thus $\partial G$ is a nonzero
Bott--Chern class.  The Dolbeault image of $\mathsf B$ is zero because
$a^\dagger=\bar\partial\bar G$, while
$\mathsf B=\overline{\mathsf A}\neq0$.  Hence
$\mathsf A,\mathsf B$ are independent, and
\cref{eq:bcpr-nonzero-dimensions} makes them a basis.

The classes $d_1\mathfrak g$ and $\mathfrak q$ are independent in
$H^{1,1}_{\bar\partial}(X)$: the first restricts to zero on every smooth
torus fibre, while the second restricts to the nonzero class
by \cref{eq:bcpr-q-restriction,eq:bcpr-q-square-integral}.  Exactness of
\cref{eq:bcpr-hodge-row-one} therefore gives
\begin{equation}
 d_1\mathfrak q=[\partial Q]_{\bar\partial}\neq0.
\label{eq:bcpr-c-nonzero}
\end{equation}
Thus $\mathsf C\neq0$, and conjugation gives
$\overline{\mathsf C}\neq0$.

Set
\begin{equation}
 v=[G\wedge Q]_{\bar\partial}
 \in H^{1,2}_{\bar\partial}(X).
\label{eq:bcpr-v-class}
\end{equation}
Its fibre restriction is the class represented by
$\Gamma\wedge Q_F$.  This class is nonzero, rather than merely represented
by a nonzero invariant form: wedging with the holomorphic one-form
$d\zeta_2$ and using \cref{eq:bcpr-q-gamma-product} gives
\begin{equation}
 \int_{F_b}Q_F\wedge\Gamma\wedge d\zeta_2
 =\frac1{4y\Delta}
   \int_{F_b}d\zeta_1\wedge d\zeta_2\wedge
                    d\bar\zeta_1\wedge d\bar\zeta_2
 =-1
\label{eq:bcpr-q-gamma-serre-pairing}
\end{equation}
for the complex orientation: the middle integral equals
$-4y\Delta>0$.  Thus the displayed constant top form has nonzero integral.
Dolbeault--Serre duality on $F_b$ therefore gives
$[\Gamma\wedge Q_F]_{\bar\partial}\ne0$.  Equivalently, if
$Q_F\wedge\Gamma=\bar\partial_F\Psi$ with $\Psi$ of bidegree $(1,1)$,
then, since $\bar\partial_Fd\zeta_2=0$,
\[
 Q_F\wedge\Gamma\wedge d\zeta_2
 =\bar\partial_F(\Psi\wedge d\zeta_2),
\]
and Stokes' theorem would make the integral in
\cref{eq:bcpr-q-gamma-serre-pairing} zero.  Here
$Q_F\wedge\Gamma$ has bidegree $(1,2)$,
$d\zeta_2$ has bidegree $(1,0)$, and
$\Psi\wedge d\zeta_2$ has bidegree $(2,1)$; its $\partial_F$ is zero in
complex dimension two, so the displayed $\bar\partial_F$-boundary is a
$d$-boundary.  Thus $v\neq0$, and it spans
$H^{1,2}_{\bar\partial}(X)$.  Injectivity in
\cref{eq:bcpr-hodge-row-two} gives
\begin{equation}
 d_1v=[\partial(G\wedge Q)]_{\bar\partial}\neq0.
\label{eq:bcpr-d-nonzero}
\end{equation}
Hence $\mathsf D\neq0$, and conjugation gives
$\overline{\mathsf D}\neq0$.  The latter has zero Dolbeault image because
its representative is $\bar\partial$-exact, so it is independent of
$\mathsf D$; \cref{eq:bcpr-nonzero-dimensions} gives the asserted
$(2,2)$ basis.

The class $[Q^2]_{\bar\partial}$ has nonzero fibre restriction by
\cref{eq:bcpr-q-square-integral}.  Since $v$ spans
$H^{1,2}_{\bar\partial}(X)$, the image of $d_1$ is spanned by $d_1v$.
On a smooth fibre,
\[
 i_b^*(d_1v)
 =d_1[\Gamma\wedge Q_F]_{\bar\partial}
 =[\partial_F(\Gamma\wedge Q_F)]_{\bar\partial}=0,
\]
because $\Gamma$ and $Q_F$ have constant fibre coefficients.  Thus every
class in $\operatorname{im}(d_1\colon H^{1,2}_{\bar\partial}\to
H^{2,2}_{\bar\partial})$ restricts to zero.
Therefore $[Q^2]_{\bar\partial}$ is independent of $d_1v$.  Exactness of
\cref{eq:bcpr-hodge-row-two} yields
\begin{equation}
 d_1[Q^2]_{\bar\partial}
 =[\partial(Q^2)]_{\bar\partial}\neq0.
\label{eq:bcpr-e-nonzero}
\end{equation}
Hence $\mathsf E\neq0$, and conjugation gives
$\overline{\mathsf E}\neq0$.  Finally, $\int_X\nu=1$ proves
$\mathsf T\neq0$.  The dimensions in
\cref{eq:bcpr-nonzero-dimensions} complete the proof.
\end{proof}

\subsection{Every Bott--Chern product}

\begin{theorem}[square-zero positive ideal]
\label{thm:bcpr-square-zero-algebra}
\Verified\ Let $N\subset H^{*,*}_{BC}(X)$ be the span of the nine positive-degree
generators in \cref{eq:bcpr-generators}.  Then
\begin{equation}
 H^{*,*}_{BC}(X)=\mathbb C\mathbf1\oplus N,
 \qquad N^2=0.
\label{eq:bcpr-algebra}
\end{equation}
Equivalently, every product of two positive-degree Bott--Chern classes is
zero.
\end{theorem}

\begin{proof}
We first calculate the products landing below top degree.  Since
$\bar\partial G=\bar\partial Q=0$ and, after conjugation,
$\partial\bar G=\partial\bar Q=0$, the graded Leibniz rule gives the exact
identities
\begin{equation}
 \partial\bar\partial(G\wedge\bar G)
 =-a\wedge a^\dagger,
\label{eq:bcpr-product-ab}
\end{equation}
\begin{equation}
 \partial\bar\partial(G\wedge\bar Q)
 =-a\wedge c^\dagger,
 \qquad
 \partial\bar\partial(\bar G\wedge Q)
 =a^\dagger\wedge c.
\label{eq:bcpr-product-acbar-bc}
\end{equation}
Consequently
\begin{equation}
 \mathsf A\mathsf B=0,\qquad
 \mathsf A\overline{\mathsf C}=0,\qquad
 \mathsf B\mathsf C=0.
\label{eq:bcpr-first-zero-products}
\end{equation}

Choose the pullback Fubini--Study metric on
$f^*\mathcal O_B(1)$.  Its Chern form is the pullback of a $(1,1)$ form on
the curve $B$, so its square vanishes pointwise.  Squaring
\cref{eq:bcpr-c1-pullback} and using
\cref{eq:bcpr-first-zero-products} gives
\begin{equation}
 \mathsf A^2+\mathsf B^2=0.
\label{eq:bcpr-sum-of-squares}
\end{equation}
\Cref{eq:bcpr-sum-of-squares} alone does not separate the two
squares.  Apply $\pi_{\bar\partial}^{2,2}$.  The class
$\mathsf B^2$ has zero Dolbeault image because
\begin{equation}
 (a^\dagger)^2
 =\bar\partial(\bar G\wedge a^\dagger).
\label{eq:bcpr-b-square-dolbeault-zero}
\end{equation}
It follows from \cref{eq:bcpr-sum-of-squares} that
$[a^2]_{\bar\partial}=0$.  On the other hand,
\begin{equation}
 [a^2]_{\bar\partial}
 =d_1[G\wedge a]_{\bar\partial},
 \qquad
 \partial(G\wedge a)=a^2.
\label{eq:bcpr-a-square-d1}
\end{equation}
The first arrow in \cref{eq:bcpr-hodge-row-two} is injective, so
\cref{eq:bcpr-a-square-d1} implies
$[G\wedge a]_{\bar\partial}=0$.  There is therefore a form
$\Phi\in\mathcal A^{1,1}(X)$ such that
\begin{equation}
 G\wedge a=\bar\partial\Phi.
\label{eq:bcpr-phi-primitive}
\end{equation}
Applying $\partial$ gives
\begin{equation}
 a^2=\partial\bar\partial\Phi,
\label{eq:bcpr-a-square-bc-zero}
\end{equation}
and hence $\mathsf A^2=0$.  Conjugation gives
$\mathsf B^2=0$.

It remains among the products below top bidegree to prove
$\mathsf A\mathsf C=0$; its conjugate will give
$\mathsf B\overline{\mathsf C}=0$.  First, the natural map
\begin{equation}
 \pi_{\bar\partial}^{3,2}\colon H^{3,2}_{BC}(X)
 \longrightarrow H^{3,2}_{\bar\partial}(X)
\label{eq:bcpr-bc-to-dolbeault-32}
\end{equation}
is an isomorphism.  Indeed, both spaces have dimension one by
\cref{eq:bcpr-nonzero-dimensions,cor:exact-hodge-diamond}, and
\cref{eq:bcpr-e-nonzero} says that $\mathsf E$ has nonzero image.
Pair $\pi_{\bar\partial}^{3,2}(\mathsf A\mathsf C)$ with the generator
$[G]_{\bar\partial}\in H^{0,1}_{\bar\partial}(X)$.  By
\cref{eq:bcpr-phi-primitive}, $\bar\partial c=0$, and Stokes' theorem,
\begin{equation}
 \int_Xa\wedge c\wedge G
 =-\int_XG\wedge a\wedge c
 =-\int_X\bar\partial\Phi\wedge c
 =-\int_X\bar\partial(\Phi\wedge c)=0.
\label{eq:bcpr-serre-pairing-zero}
\end{equation}
Dolbeault--Serre duality makes
\begin{equation}
 H^{3,2}_{\bar\partial}(X)\times
 H^{0,1}_{\bar\partial}(X)\longrightarrow\mathbb C,
 \qquad
 ([\alpha],[\rho])\longmapsto\int_X\alpha\wedge\rho,
\label{eq:bcpr-serre-pairing}
\end{equation}
a perfect pairing.  Since
$H^{0,1}_{\bar\partial}(X)=\mathbb C[G]_{\bar\partial}$,
\cref{eq:bcpr-serre-pairing-zero} implies
$\pi_{\bar\partial}^{3,2}(\mathsf A\mathsf C)=0$; injectivity in
\cref{eq:bcpr-bc-to-dolbeault-32} gives
\begin{equation}
 \mathsf A\mathsf C=0,
 \qquad
 \mathsf B\overline{\mathsf C}=0.
\label{eq:bcpr-second-zero-products}
\end{equation}
This step uses Dolbeault--Serre duality together with the independently
proved injectivity of \cref{eq:bcpr-bc-to-dolbeault-32}.  The
Bott--Chern--Aeppli pairing of
\cite[\S2.c]{Schweitzer2007} by itself would supply a perfect pairing but
would not annihilate this product.

The only products not yet considered which can have bidegree at most
$(3,3)$ are
\begin{equation}
 \mathsf A\mathsf D,\quad
 \mathsf A\overline{\mathsf D},\quad
 \mathsf B\mathsf D,\quad
 \mathsf B\overline{\mathsf D},\quad
 \mathsf C\overline{\mathsf C},
\label{eq:bcpr-top-products}
\end{equation}
all of bidegree $(3,3)$.  Every non-top generator in
\cref{eq:bcpr-generators} is de Rham exact:
\begin{equation}
\begin{gathered}
 a=dG,\qquad a^\dagger=d\bar G,\qquad
 c=dQ,\qquad c^\dagger=d\bar Q,\\
 d_0=d(G\wedge Q),\qquad
 d_0^\dagger=d(\bar G\wedge\bar Q),\qquad
 e=d(Q^2),\qquad e^\dagger=d(\bar Q^2).
\end{gathered}
\label{eq:bcpr-de-rham-primitives}
\end{equation}
Each factor paired with an exact factor is $d$-closed, so every product in
\cref{eq:bcpr-top-products} is $d$-exact and has integral zero.  Integration
induces an isomorphism
\begin{equation}
 H^{3,3}_{BC}(X)\xrightarrow{\ \int_X\ }\mathbb C.
\label{eq:bcpr-top-integration}
\end{equation}
It is well defined by Stokes' theorem, the source has dimension one by
\cref{eq:bcpr-nonzero-dimensions}, and it is nonzero on $\mathsf T$.
Therefore all five products in \cref{eq:bcpr-top-products} vanish in
Bott--Chern cohomology.

All other products of positive-degree generators have $p>3$ or $q>3$ and
vanish already as differential forms.
\Cref{eq:bcpr-first-zero-products,eq:bcpr-a-square-bc-zero,eq:bcpr-second-zero-products,eq:bcpr-top-products,eq:bcpr-top-integration}
exhaust every possible positive-degree product
and prove \cref{eq:bcpr-algebra}.
\end{proof}

\begin{remark}[literal global gauge boundary]
\label{rem:bcpr-global-gauge-boundary}
The Bott--Chern algebra \cref{eq:bcpr-algebra} and the fibre-coordinate
restrictions displayed in
\cref{eq:bcpr-gamma-form,eq:bcpr-u-form,eq:bcpr-q-fibre,eq:bcpr-q-fibre-expanded,eq:bcpr-q-gamma-product,eq:bcpr-q-square}
are complete.  The canonical data do not select a preferred literal global
coefficient atlas for the smooth forms $G$ and $Q$.  Such an atlas would
require explicit smooth local trivializations and horizontal correction
functions through both finite affine fillings and the cusp.  After fixing
$\mathfrak g$ and one coherent-Leray lift $\mathfrak q$, changing those
smooth gauges changes $G$ and $Q$ by Dolbeault coboundaries but changes
none of the cohomology classes, nonvanishing arguments, or products proved
above.  This is a limitation only on a preferred global coefficient
presentation, not on the representatives as cohomology classes or on the
Bott--Chern algebra.
\end{remark}
